\documentclass[11pt]{article}
\usepackage[margin=0.75in]{geometry}
\usepackage{amsmath,amssymb,siunitx}
\usepackage{parskip}
\usepackage{enumitem}
\usepackage{multicol}

\title{PS160 --- Master Equation Sheet\\\large Fluids, Oscillations, Waves, Sound, Thermodynamics, EM Waves, Optics}
\author{Full Course (Modules 12--20, 33--36)}
\date{}

\begin{document}
\maketitle

\noindent This is a single-document reference for every equation used in PS160. It is organized by exam and by module. See the individual \texttt{equations.tex} files in each \texttt{midterm\_0X/} folder for topic-focused cheat sheets.

\paragraph{Note on the official Final equation sheet.}
The official Final exam equation sheet (see
\texttt{midterm\_03/test\_finalexam\_DRAFT\_ps160\_2024\_fall.pdf}, pp.\ 2--3)
covers only Midterm~1 and Midterm~2 material (Modules 12, 14--20).
\textbf{No Module 33--36 formulas are provided.} Items below tagged
\fbox{\textsc{sheet}} are on the official sheet; items tagged
\fbox{\textsc{memorize}} must be memorized or derived.
See \texttt{midterm\_03/equation\_gaps.md} for the full audit.

\tableofcontents
\newpage

%==========================================================
\part*{Midterm 1 --- Mechanics of Fluids and Waves}
%==========================================================

\section{Fluid Mechanics (Module 12)}
\fbox{\textsc{sheet}} --- everything in this section is on the official sheet.

\[
\rho = \frac{m}{V}\qquad p = \frac{F_\perp}{A}
\]
\[
\frac{dp}{dy} = -\rho g\qquad p(y) = p_0 - \rho g y
\]
\paragraph{Buoyancy (Archimedes).} $F_B = \rho_{\text{fluid}} V_{\text{disp}} g$.
\paragraph{Continuity.} $\rho A v = \text{const}$ (incompressible: $A_1 v_1 = A_2 v_2$).
\paragraph{Bernoulli.} $p + \rho g y + \tfrac{1}{2}\rho v^{2} = \text{const}$.
\paragraph{Torricelli.} $v = \sqrt{2 g h}$.
\paragraph{Volume flow rate.} $dV/dt = A v$.

\section{Oscillations (Module 14)}
SHM equation, $x(t)$, three $\omega$ forms, and the $f$/$T$/$\omega$ relations
are \fbox{\textsc{sheet}}. Energy, damped, and driven results are
\fbox{\textsc{memorize}}.

\paragraph{SHM equation.} $\ddot{x} + \omega^{2}x = 0$, solution $x(t) = A\cos(\omega t + \varphi)$.
\paragraph{Period relations.} $\omega = 2\pi f = 2\pi/T$, $T = 2\pi/\omega$.
\paragraph{Angular frequencies.}
\[
\omega_{\text{spring}} = \sqrt{k/m},\quad
\omega_{\text{simple pend.}} = \sqrt{g/\ell},\quad
\omega_{\text{phys. pend.}} = \sqrt{mgd/I}
\]
\paragraph{Energy.} $E = \tfrac{1}{2} k A^{2} = \tfrac{1}{2} m\omega^{2} A^{2}$.
\paragraph{Damped.} $m\ddot x + b\dot x + kx = 0$, $\omega' = \sqrt{k/m - b^{2}/(4m^{2})}$.
\paragraph{Driven, steady state amplitude.}
$A_d = F_0/\sqrt{(k - m\omega_d^{2})^{2} + b^{2}\omega_d^{2}}$.

\section{Mechanical Waves (Module 15)}
Wave function, $k$, $v=\lambda f$, all three speed formulas, $f_n$, $\lambda_n$,
and $P_{\text{av}}$ are \fbox{\textsc{sheet}}. The wave equation itself is
\fbox{\textsc{memorize}}.

\paragraph{Wave function.} $y(x,t) = A\cos(kx - \omega t)$, $k = 2\pi/\lambda$, $v = \lambda f = \omega/k$.
\paragraph{Wave equation.} $\partial_t^{2} y = v^{2}\partial_x^{2} y$.
\paragraph{Speeds.}
\[
v_{\text{string}} = \sqrt{F/\mu},\quad
v_{\text{solid}} = \sqrt{Y/\rho},\quad
v_{\text{fluid}} = \sqrt{B/\rho}
\]
\paragraph{Power (string).} $P_{\text{av}} = \tfrac{1}{2}\sqrt{\mu F}\,\omega^{2}A^{2}$.
\paragraph{Standing waves.}
\[
f_n = \frac{nv}{2L}\;\text{or}\;\frac{nv}{4L},\qquad
\lambda_n = \frac{2L}{n}\;\text{or}\;\frac{4L}{n}
\]

\section{Sound (Module 16)}
Intensity, $p_{\max}$, spherical-source $I = P/(4\pi r^2)$, decibels, Doppler,
and beats are \fbox{\textsc{sheet}}. Temperature dependence of $v$ is
\fbox{\textsc{memorize}}.

\paragraph{Sound speed.} $v = \sqrt{B/\rho}$; air $\approx \SI{343}{m/s}$.
\paragraph{Temperature dependence (air) \fbox{\textsc{memorize}}.} $v(T) = (331\;\mathrm{m/s})\sqrt{T/(273\;\mathrm{K})}$.
\paragraph{Intensity.} $I = \tfrac{1}{2}\sqrt{\rho B}\,\omega^{2}A^{2} = p_{\max}^{2}/(2\sqrt{\rho B})$.
\paragraph{Pressure amplitude.} $p_{\max} = BkA$.
\paragraph{Spherical source.} $I = P/(4\pi r^{2})$.
\paragraph{Decibels.} $\beta = (10\text{ dB})\log_{10}(I/I_0)$, $I_0 = \SI{e-12}{W/m^{2}}$.
\paragraph{Doppler.} $f_L = f_S\,(v + v_L)/(v + v_S)$ (source right of listener).
\paragraph{Beats.} $f_{\text{beat}} = |f_a - f_b|$.

%==========================================================
\part*{Midterm 2 --- Thermodynamics}
%==========================================================

\section{Temperature \& Heat (Module 17)}
Expansion, $Q = mc\Delta T$, latent heat, conduction, and radiation $H = A\sigma T^4$
are \fbox{\textsc{sheet}}. Temperature conversions and thermal stress are
\fbox{\textsc{memorize}}. The sheet's radiation form omits the emissivity $e$;
include it from memory if a problem specifies $e \ne 1$.

\paragraph{Temperature scales \fbox{\textsc{memorize}}.} $T_K = T_C + 273.15$; $T_F = \tfrac{9}{5}T_C + 32$; $T_R = T_F + 459.67$.
\paragraph{Expansion.}
$\Delta L = \alpha L_0 \Delta T$, $\Delta V = \beta V_0\Delta T$, $\beta = 3\alpha$.
\paragraph{Thermal stress \fbox{\textsc{memorize}}.} $F/A = -Y\alpha\Delta T$.
\paragraph{Heat.} $Q = mc\Delta T = nC\Delta T$.
\paragraph{Latent heat.} $Q = \pm m L$.
\paragraph{Conduction.} $H = kA\Delta T/L$.
\paragraph{Radiation.} Sheet: $H = A\sigma T^{4}$. With emissivity (\fbox{\textsc{memorize}}): $H = eA\sigma T^4$. $\sigma = \SI{5.67e-8}{W/(m^{2} K^{4})}$.

\section{Kinetic Theory (Module 18)}
Ideal gas law, equipartition, $v_{rms}$, MB distribution, $C_V/C_P/\gamma$, and
$C_V=3R$ for solids are \fbox{\textsc{sheet}}. Total translational KE form is
derivable.

\paragraph{Ideal gas.} $pV = NkT = nRT$, $kN_A = R$.
\paragraph{Equipartition.} $\tfrac{1}{2}m\langle v^{2}\rangle = \tfrac{3}{2}kT$.
\paragraph{Total translational KE.} $K_{\text{tr,tot}} = \tfrac{3}{2}nRT$ for $n$ moles.
\paragraph{rms speed.} $v_{\text{rms}} = \sqrt{3kT/m}$.
\paragraph{Maxwell--Boltzmann.} $f(v)\propto v^{2}e^{-mv^{2}/(2kT)}$.
\paragraph{Molar heat capacities.}
$C_V = (f/2)R$, $C_P = C_V + R$, $\gamma = C_P/C_V$; solid: $C_V = 3R$.

\section{First Law (Module 19)}
$dU = dQ - dW$, $W = \int p\,dV$, and $W = Q = Q_H + Q_C$ for a cycle are
\fbox{\textsc{sheet}}. The closed-form process results below are
\fbox{\textsc{memorize}} (derive from $W = \int p\,dV$ if needed).

\paragraph{First law.} $dU = \bar dQ - \bar dW$, $W = \int p\,dV$.
\paragraph{Ideal gas internal energy.} $U = nC_V T$, $\Delta U = nC_V\Delta T$.

\paragraph{Processes (ideal gas) \fbox{\textsc{memorize}}.}
\begin{itemize}\itemsep0pt
\item Isobaric: $W = p\Delta V$, $Q = nC_P\Delta T$.
\item Isochoric: $W = 0$, $Q = nC_V\Delta T = \Delta U$.
\item Isothermal: $\Delta U = 0$, $Q = W = nRT\ln(V_2/V_1)$.
\item Adiabatic: $Q = 0$, $pV^{\gamma} = \text{const}$, $TV^{\gamma-1} = \text{const}$, $W = (p_1V_1 - p_2V_2)/(\gamma - 1)$.
\end{itemize}

\section{Second Law, Entropy, Cycles (Module 20)}
Zeroth law, $\Delta S \ge 0$, $\Delta S = \int dQ/T$, $S = k\ln w$, $e = W/Q_H$,
$K = Q_C/(-W)$, $K = -Q_H/(-W)$ for a heat pump in heating mode, and the
Carnot entropy identity $\Delta S = Q_H/T_H + Q_C/T_C = 0$ are all
\fbox{\textsc{sheet}}. The Carnot efficiency form $e = 1 - T_C/T_H$ is
\fbox{\textsc{memorize}} (derive by combining $e = W/Q_H = (Q_H+Q_C)/Q_H$ with
$\Delta S = 0$).

\paragraph{Zeroth law.} $T_A = T_B = T_C$.
\paragraph{Second law.} $\Delta S_{\text{universe}} \ge 0$.
\paragraph{Engine cycle.} $W_{\text{net}} = Q_H + Q_C$; efficiency $e = W/Q_H = 1 - |Q_C|/Q_H$.
\paragraph{Carnot \fbox{\textsc{memorize}}.} $e_{\text{Carnot}} = 1 - T_C/T_H$.
\paragraph{Refrigerator (sheet form).} $K = Q_C/(-W)$.
\paragraph{Heat pump, heating mode (sheet form).} $K = -Q_H/(-W)$.
\paragraph{Entropy.} $\Delta S = \int dQ_{\text{rev}}/T$, $S = k\ln w$.
\paragraph{Carnot entropy.} $\Delta S = Q_H/T_H + Q_C/T_C = 0$.
\paragraph{Useful simple processes \fbox{\textsc{memorize}} (derive from $\Delta S = \int dQ/T$).}
\[
\Delta S_{\text{isoT, ideal gas}} = nR\ln\frac{V_2}{V_1},\quad
\Delta S_{\text{heating}} = mc\ln\frac{T_2}{T_1},\quad
\Delta S_{\text{phase change}} = \pm\frac{mL}{T}
\]

%==========================================================
\part*{Midterm 3 --- Optics and Electromagnetic Waves}
%==========================================================

\section{EM Waves (Module 33)}
\fbox{\textsc{memorize}} --- \emph{nothing} from this section is on the official
sheet. Memorize or be ready to derive every formula below.

\paragraph{Maxwell's equations.}
\[
\oint\vec E\cdot d\vec A = \frac{Q_{\text{enc}}}{\varepsilon_0},\quad
\oint\vec B\cdot d\vec A = 0
\]
\[
\oint\vec E\cdot d\vec\ell = -\frac{d\Phi_B}{dt},\quad
\oint\vec B\cdot d\vec\ell = \mu_0\!\left(I_{\text{enc}} + \varepsilon_0\frac{d\Phi_E}{dt}\right)
\]
\paragraph{Speed of light.} $c = 1/\sqrt{\mu_0\varepsilon_0} \approx \SI{3.00e8}{m/s}$.
\paragraph{Plane wave relations.} $E/B = c$, $\vec E\perp\vec B\perp\vec k$, $v = \lambda f = \omega/k$.
\paragraph{In a medium.} $v = c/n$, $\lambda_{\text{med}} = \lambda_0/n$.
\paragraph{Energy density.} $u = \varepsilon_0 E^{2} = B^{2}/\mu_0$.
\paragraph{Poynting vector / intensity.}
$\vec S = (1/\mu_0)\vec E\times\vec B$,
$I = \tfrac{1}{2}c\varepsilon_0 E_{\max}^{2} = E_{\max}^{2}/(2\mu_0 c)$.
\paragraph{Radiation pressure.} $p_{\text{abs}} = I/c$, $p_{\text{refl}} = 2I/c$.
\paragraph{Malus's law.} $I = I_{\max}\cos^{2}\phi$; unpolarized through ideal polarizer $\to I_0/2$.
\paragraph{Brewster.} $\tan\theta_B = n_2/n_1$.

\section{Geometric Optics (Module 34)}
\fbox{\textsc{memorize}} --- nothing from this section is on the official sheet.

\paragraph{Reflection.} $\theta_i = \theta_r$.
\paragraph{Snell.} $n_1\sin\theta_1 = n_2\sin\theta_2$.
\paragraph{Index.} $n = c/v$, $\lambda_{\text{med}} = \lambda_0/n$.
\paragraph{Total internal reflection.} $\sin\theta_c = n_2/n_1$ (with $n_1 > n_2$).
\paragraph{Mirror / thin lens equation.}
$1/s + 1/s' = 1/f$, $f_{\text{mirror}} = R/2$.
\paragraph{Magnification.} $m = -s'/s$.
\paragraph{Lensmaker.} $1/f = (n - 1)(1/R_1 - 1/R_2)$.
\paragraph{Single refracting surface.} $n_1/s + n_2/s' = (n_2 - n_1)/R$.
\paragraph{Angular magnification.}
$M_{\text{magnifier}} = (25\text{ cm})/f$, $M_{\text{telescope}} = -f_{\text{obj}}/f_{\text{eye}}$.
\paragraph{Lens power.} $P\,[\text{diopter}] = 1/f\,[\text{m}]$. Converging $P>0$; diverging $P<0$.
\paragraph{Vision correction.} Myopia (near-sightedness) $\to$ diverging lens; hyperopia (far-sightedness) $\to$ converging lens.
\paragraph{Apparent depth.} $d_{\text{app}} = d_{\text{real}}(n_{\text{obs}}/n_{\text{med}})$.

\section{Interference (Module 35)}
\fbox{\textsc{memorize}} --- nothing from this section is on the official sheet.

\paragraph{Path-difference phases.} $\Delta\phi = (2\pi/\lambda)\Delta r$.
\paragraph{Two-source conditions.}
$\Delta r = m\lambda$ (bright), $\Delta r = (m + \tfrac{1}{2})\lambda$ (dark).
\paragraph{Young's double slit.}
$d\sin\theta = m\lambda$, $y_m \approx m\lambda R/d$, fringe spacing $\lambda R/d$.
\paragraph{Central max width (double slit).} $w_{\text{central}} = 2\lambda R/d$ (between $m=\pm 1$ dark fringes on either side).
\paragraph{Two-source intensity.} $I(\theta) = I_0\cos^{2}[(\pi d\sin\theta)/\lambda]$.
\paragraph{Thin film (normal incidence).}
$2 n t = m\lambda_0$ or $(m + \tfrac{1}{2})\lambda_0$ (swap if exactly one hard reflection).
\paragraph{Michelson interferometer.} $2\Delta d = \Delta m\,\lambda$.

\section{Diffraction (Module 36)}
\fbox{\textsc{memorize}} --- nothing from this section is on the official sheet.

\paragraph{Single-slit minima.} $a\sin\theta = m\lambda$, $m = \pm 1,\pm 2,\dots$
\paragraph{Central max width (single slit).} $w_{\text{central}} = 2\lambda R/a$ (between $m=\pm 1$ dark fringes).
\paragraph{Grating visible orders.} $m_{\max} = \lfloor d/\lambda\rfloor$; total visible maxima $=2m_{\max}+1$.
\paragraph{Single-slit intensity.}
$I = I_0\,[\sin(\beta/2)/(\beta/2)]^{2}$, $\beta = 2\pi a\sin\theta/\lambda$.
\paragraph{Double-slit with finite width (combined).}
$I = I_0\cos^{2}(\pi d\sin\theta/\lambda)\cdot[\sin(\beta/2)/(\beta/2)]^{2}$.
\paragraph{Grating maxima.} $d\sin\theta = m\lambda$; resolving power $R = \lambda/\Delta\lambda = Nm$.
\paragraph{Circular aperture / Rayleigh.} $\sin\theta_R = 1.22\lambda/D$.
\paragraph{Bragg.} $2 d\sin\theta = m\lambda$ ($\theta$ measured from the planes).

%==========================================================
\section*{Physical Constants}
%==========================================================
\begin{tabular}{ll}
Speed of light in vacuum & $c = \SI{2.998e8}{m/s}$ \\
Gravitational acceleration & $g = \SI{9.81}{m/s^{2}}$ \\
Density of water & $\rho_{\text{water}} = \SI{1000}{kg/m^{3}}$ \\
Atmospheric pressure & $p_{\text{atm}} = \SI{101325}{Pa}$ \\
Specific heat of water & $c_{\text{water}} = \SI{4190}{J/(kg\cdot K)}$ \\
Boltzmann constant & $k = \SI{1.380649e-23}{J/K}$ \\
Avogadro's number & $N_A = \SI{6.022e23}{mol^{-1}}$ \\
Universal gas constant & $R = \SI{8.314}{J/(mol\cdot K)}$ \\
Stefan--Boltzmann constant & $\sigma = \SI{5.670e-8}{W/(m^{2}K^{4})}$ \\
Permittivity of free space & $\varepsilon_0 = \SI{8.854e-12}{F/m}$ \\
Permeability of free space & $\mu_0 = 4\pi\times\SI{e-7}{T\cdot m/A}$ \\
Sound speed in air & $\approx \SI{343}{m/s}$ \\
\end{tabular}

\end{document}
